\section{Approaches We Considered}

This section follows the project chronologically: what we tried, what we measured, and why each idea was kept or dropped.

Our investigative trajectory progressed chronologically from initial architectural deliberations through empirical dead ends, leading to the development of a zero-parameter consensus gate, its noise-adaptive generalization, and an online noise-routing subsystem.

\subsection{Architectural Conception and Multi-AI Deliberation}

The project originated from the hierarchical decoding paradigm formalized in recent literature (e.g., arXiv:2611.16271), wherein a rapid, low-complexity filter resolves typical, low-weight syndrome configurations, delegating only recalcitrant syndromes to a heavy, computationally intensive backend. 

To explore this design space, our team established a structured tri-agent collaborative workflow comprising Claude (operating via Claude Code), GPT (via Codex), and Gemini (via agy). The initial debate focused on the granularity of delegation:
\begin{itemize}
    \item \textit{Partial-Shot Commitment vs.\ Whole-Shot Escalation}: One proposed approach was to run a lightweight Belief Propagation (BP) decoder, identify high-marginal ``confident'' data bits, freeze their values, and pass the reduced Tanner graph to an Ordered Statistics Decoder (BP-OSD) or Relay-BP. All three agents reached a consensus to reject this mechanism. In bivariate bicycle codes, check operators possess weight 6 and wrap globally around the $12 \times 6$ torus. Freezing bits based on local marginals risks locking the system into an incorrect topological homology (a spurious logical coset), guaranteeing a logical failure that no downstream decoder can overturn. Consequently, the team restricted all designs to whole-shot accept/escalate architectures.
    \item \textit{Evaluation Metrics}: The panel agreed that intermediary metrics, such as router classification accuracy, were misleading. System performance had to be evaluated strictly on end-to-end paired logical error rate per cycle, defined as $\text{LER} = 1 - (1 - P_{\text{pair\_fail}})^{1/12}$ across coupled $X$ and $Z$ memory shots, alongside mean and 99th-percentile (p99) execution latency.
    \item \textit{Scope Delimitation}: Although the challenge brief highlighted qubit loss and erasure heralding as a bonus topic, the provided simulator kit did not model loss channels. The panel deferred erasure decoding to prioritize the circuit-level depolarizing sweep.
\end{itemize}

\subsection{Early Empirical Dead Ends at $p = 0.004$}

Early calibration experiments demonstrated that a pure 80-iteration min-sum BP stage (BP-80) converged cleanly on $87\%$, $64\%$, and $30\%$ of shots at $p = 0.003$, $0.004$, and $0.005$, respectively, executing in only 5--9~ms. Crucially, across initial 500-shot samples, a converged BP-80 decode was never incorrect (0/437, 0/318, and 0/148 failures). However, constructing a practical filter around BP proved challenging:

\begin{itemize}
    \item \textit{Gemini's Low-Weight Acceptance Heuristic}: Gemini initially hypothesized that BP solutions should be accepted if and only if the total predicted correction weight satisfied $|e| \le 4$. Empirical validation at $p = 0.004$ immediately falsified this hypothesis: the median weight of true, syndrome-consistent corrections across 12 cycles of noise is 45. A threshold of 4 was far below the weight of typical correct corrections.
    \item \textit{Direct BP Pre-Gating}: Placing BP-80 directly in front of Relay-BP yielded marginal latency improvements, shifting mean runtime from 266~ms to 254~ms. Because BP-80 fails on approximately $36\%$ of shots at $p = 0.004$, the overall system latency remained dominated by the intractable tail of unresolved shots.
    \item \textit{BP-LSD as an Intermediate Tier}: Localized Statistics Decoding (BP-LSD) was evaluated as a mid-tier alternative to full OSD. At $p = 0.004$, BP-LSD produced a cycle LER of $9.5 \times 10^{-3}$, substantially worse than standard BP-OSD ($4.4 \times 10^{-3}$), and was dropped.
    \item \textit{Single-Run Confidence and Machine-Learned Gating}: Attempts to gate a single run of fast Relay-BP ($\text{stop\_nconv}=1$) using internal solver diagnostics (iterations used, correction weight, log-likelihood ratio margins) via logistic regression and Gradient Boosted Decision Trees (GBDT) proved counterproductive. The best heuristic filter reduced runtime by only $14\%$ while increasing failure counts from 12 to 13; the GBDT filter produced 5 additional logical failures over baseline Relay-BP.
    \item \textit{Oracle Headroom Assessment}: To ascertain whether tiered gating was fundamentally viable, we simulated an oracle gate possessing perfect prescience of whether a fast decode matched the true logical class. The oracle achieved 8 paired failures at 128~ms mean latency, compared to Relay-BP's 12 failures at 361~ms. This confirmed substantial theoretical headroom, provided a sufficiently reliable gating criterion could be devised.
\end{itemize}

\subsection{The Agreement Gate and Lean Execution}

The breakthrough emerged from exploiting the stochastic exploration dynamics of Relay-BP itself. Standard Relay-BP (Mueller et al., relay\_bp 0.2.2) executes an initial 80-iteration BP leg, followed by successive 60-iteration legs perturbed by random memory coefficients $\gamma \in (-0.226, 0.622)$. It accumulates solutions until finding $\text{stop\_nconv} = 5$ distinct valid candidates (capped at 300 legs) and returns the lowest-energy configuration.

Rather than relying on learned classifiers, we introduced the \textit{Agreement Gate}:
\begin{enumerate}
    \item Execute a truncated Relay-BP instance ($\text{stop\_nconv} = 1$, maximum 50 legs) initialized with seed~0.
    \item Execute a second truncated Relay-BP instance ($\text{stop\_nconv} = 1$, maximum 50 legs) initialized with seed~1.
    \item If both instances converge to identical 12-bit logical predictions, the shot is accepted immediately.
    \item If either instance fails to find a valid syndrome solution, or if their 12-bit logical predictions disagree, the shot is escalated to the full ``referee'' decoder (standard kit Relay-BP with $\text{stop\_nconv} = 5$, 300 legs).
\end{enumerate}

Because the first 80 iterations of Relay-BP are completely deterministic (free of random memory perturbations), seeds~0 and 1 trace identical execution trajectories during this initial leg. This insight enabled the \textit{Lean Agreement Gate}:
\begin{itemize}
    \item If seed~0 converges within the initial 80 iterations, accept immediately. Seed~1 would behave identically; across approximately 34,000 empirical test shots, zero discrepancies were observed on this path.
    \item If seed~0 finds no solution at all within its leg cap, escalate immediately to avoid redundant computation.
    \item Otherwise (if seed~0 required perturbed legs to converge), run seed~1. If seed~1 agrees with seed~0, accept; otherwise, escalate to the referee.
\end{itemize}

\subsection{Pre-Registered Verification}

Under a pre-registered protocol (Document~05) evaluating 17,000 fresh paired shots, the Agreement Gate matched the fidelity of full Relay-BP across the board (pooled failures: 317 for the gate vs.\ 327 for Relay-BP). At $p = 0.005$, the Agreement Gate demonstrated a statistically significant accuracy enhancement over baseline Relay-BP, recording 14 favorable discordant pairs against only 2 unfavorable pairs (McNemar test $p \approx 0.004$).

In terms of latency, the original agreement gate achieved mean runtime ratios of $0.48\times$ at $p = 0.003$ and $0.70\times$ at $p = 0.004$ ($0.79\times$ in single-worker controlled tests). However, at $p = 0.005$, mean latency was $0.96\times$ (failing the pre-registered acceleration criterion), while p99 latency degraded by $13\text{--}30\%$ across $p \ge 0.004$. The Lean variant resolved the mean latency bottleneck, achieving runtime ratios of $0.36\times$, $0.63\times$, and $0.88\times$ across $p \in \{0.003, 0.004, 0.005\}$ ($0.72\times$ single-worker at $0.004$), preserving full logical fidelity.

\subsection{Noise-Adaptive Parameter Optimization}

To address latency degradation at higher physical error rates, we investigated noise-adaptive configurations. The operating envelope was parameterized by the fast restart cap $K$ (limiting fast legs to $80 + 60K$ iterations) and the referee tier: kit Relay-BP ($\text{R5}$, $\text{stop\_nconv} = 5$), $\text{R2}$ ($\text{stop\_nconv} = 2$), or $\text{R1}$ ($\text{stop\_nconv} = 1$, 300 legs).

\textit{Discovery Stage}: Twelve configurations were evaluated across six noise rates ($p \in [0.003, 0.006]$) over 24,500 paired shots. Selecting the fastest configuration whose failure ratio remained within $1.05\times$ of baseline Relay-BP yielded the initial discovery schedule: $K5/\text{R2}$ at $p = 0.003$; $K50/\text{R5}$ at $p \in [0.004, 0.005]$; $K5/\text{R2}$ at $p = 0.0055$; and $K5/\text{R1}$ at $p = 0.006$. Discovery data revealed that while lightweight referees ($\text{R2}$, $\text{R1}$) degraded accuracy by $15\text{--}39\%$ at $p \in [0.004, 0.005]$, they introduced no penalty at $p \ge 0.0055$, where the escalated shots failed whichever of the tested referees handled them.

\textit{Confirmation Stage and Low-Noise Ablation}: Confirmation runs on fresh seeds exposed critical edge cases:
\begin{itemize}
    \item At $p = 0.003$, the $K5/\text{R2}$ configuration failed confirmation (9 failures vs.\ 2 for Relay-BP), isolated directly to the weaker $\text{R2}$ referee.
    \item Subsequent ablations at $p = 0.003$ using the robust $\text{R5}$ referee demonstrated that $K5$ yielded 15 vs.\ 14 failures, while $K50$ produced 10 vs.\ 7. These fluctuations demonstrated that 10,000 shots at $p = 0.003$ are statistically underpowered; Relay-BP itself exhibited 2, 14, and 7 failures across three independent 10,000-shot runs.
    \item At $p = 0.0045$ and $0.005$ accuracy passed and at $p = 0.004$ it was inconclusive (69 vs 61), while p99 latency was about $15\%$ worse than Relay-BP at all three, failing the latency criterion.
\end{itemize}

In response, we established a \textit{Revised Operational Table}: retain the full $\text{R5}$ referee with $K = 10$ across all standard operating regimes ($p \le 0.005$), deploying $K5/\text{R2}$ only at $p = 0.0055$ and $K5/\text{R1}$ at $p = 0.006$.

\subsection{Syndrome-Level Noise Estimation and Online Routing}

To deploy the noise-adaptive table without external parameter injection, we implemented an automated noise-estimation front-end to infer $p$ directly from syndrome alarms:
\begin{itemize}
    \item \textit{Detector Characteristics}: The mean count of $X$-basis detector alarms scales monotonically with circuit noise: 34, 66, 95, 122, 147, 170, and 191 alarms at $p = 0.001$ to $0.007$ in steps of $0.001$, respectively.
    \item \textit{Model Training and Evaluation}: Across a dataset of 61,500 paired shots spanning 41 noise levels in $[0.002, 0.007]$, models were trained on 24 spatiotemporal features per basis (alarm counts partitioned by cycle and check type). Over a sliding window of the last 10 pairs, a simple scalar alarm count achieved a relative estimation error of $2\text{--}3\%$, matching a multi-layer perceptron (MLP with layers 32 and 16, $2\text{--}3\%$ error) and ridge regression ($2\text{--}4\%$ error). On single shots, error was $6\text{--}9\%$, dropping to $\sim 0.8\%$ over 100 shots.
    \item \textit{Computational Efficiency and Sufficiency}: Total detector counting requires $1\ \mu\text{s}$, compared to $\sim 70\ \mu\text{s}$ per single inference for the MLP (under $1\ \mu\text{s}$ batched). Under spatially uniform depolarizing noise, the total alarm count functions as an empirical sufficient statistic; the tested ML estimators gave no measurable improvement.
    \item \textit{Dynamic Drift Streaming}: In a synthetic drift stream where $p$ shifted every 300 pairs, an alarm-counting window of $W = 10$ pairs delivered 203 total logical failures, matching both the static oracle (207 failures) and pure Relay-BP (207 failures), while reducing runtime to $0.67\times$.
    \item \textit{Asymmetric Risk Profile}: Analysis identified an asymmetric penalty landscape: under-estimating noise retains conservative, safe referee settings ($\text{R5}$). Conversely, a $+10\%$ upward bias at $p \in [0.0045, 0.005]$ triggered premature transitions to aggressive settings, inducing an $11\text{--}18\%$ surge in logical failures.
\end{itemize}
